\section{Related Work}
\label{sec:related}

\citet{wang2026prefillawareness} show that frontier LLMs exhibit prefill awareness, to differing extents. Claude Opus 4.5 for instance recognizes prefill in 35\% of positive samples (with 0\% false positives). These tests were conducted on models much larger than the ones we tested, and the authors mostly considered settings in which a substantial part of the text was prefilled.

\citet{g2026enaction} introduces the hypothesis that on-policy and off-policy generations are separated by the model using a surprisal signal of the previous token, that mechanistically leads to a higher entropy in the distribution of the next token. In their paper, they dub this increased entropy reaction "implicit recognition", as opposed to "explicit recognition" which refers to the model's ability to verbally recognize prefill when asked. They find that the explicit prefill detection is routed through a different mechanism. We also find a correlation between the signal the fine-tuned model uses in the no-thinking settings and the surprisal of the current token. This paper was one of the inspirations for us to use a surprisal-driven signal in our training (See \ref{sec:method}).

\citet{nguyen2026selfreport} study different fine-tuning methods in order to see if they are able to teach models to better recognize adversarial prefill (details in Appendix H in their paper). While they consider longer spans of tokens being prefilled from a different model, we focus in this work on recognizing a weaker signal: 1--4 planted tokens, either scattered across the previous turn or in a consecutive span (see \ref{sec:method}). In their paper, they observe that their fine-tuning didn't transfer to better performance in prefill recognition when using a ``tampering probe'' question ``Did anyone tamper with your response?''. We similarly found that our fine-tuning didn't transfer to better performance with the same probe.